\section{Detailed mathematical derivations}
\label{app:derivations}

This appendix contains the calculations that are used in the article but whose
length would interrupt it: the element matrices of the Legendre basis, five
worked matrix elements, the construction of the blocks \eqref{eq:block} from the
weak equations, the edge-matching system, the normalisation of the modal power,
and one end-to-end numerical evaluation of the reported share.

\subsection{Element matrices of the Legendre basis}

Three coordinates appear in the formulation and they must be separated before
any integral is written, because the orthonormality of the basis and the size of
the element matrices both depend on which one is meant.  The physical coordinate
is $z\in[-d/2,d/2]$.  The non-dimensional coordinate of the field equations is
$\zeta=2z/d\in[-1,1]$, with $\dd z=(d/2)\dd\zeta$.  The basis is generated on a
third, the half-interval variable
\begin{equation}
\label{eq:app_measures}
s=\frac{z}{d}=\frac{\zeta}{2}\in\Bigl[-\tfrac12,\tfrac12\Bigr],
\qquad
\dd s=\frac{\dd z}{d}=\frac{\dd\zeta}{2},
\end{equation}
and every element matrix in this appendix is an integral over $s$ with the
measure $\dd s$.  The basis is
\begin{equation}
\label{eq:app_basis}
b_j(s)=\sqrt{2j+1}\,P_j(2s),\qquad j=0,\dots,N-1,
\end{equation}
with $N$ the number of basis functions per field.  Since
$\int_{-1}^{1}P_iP_j\dd\zeta=2\delta_{ij}/(2j+1)$, the Gram matrix in the $s$
measure is
\begin{equation}
\label{eq:app_mass}
\int_{-1/2}^{1/2}b_ib_j\,\dd s
=\frac{2j+1}{2}\int_{-1}^{1}P_iP_j\,\dd\zeta=\delta_{ij},
\end{equation}
the first factor being the Jacobian of $s=\zeta/2$ and the second the Legendre
normalisation.  The same integral in the $\zeta$ measure is
$\int_{-1}^{1}b_ib_j\dd\zeta=2\delta_{ij}$: the two differ by exactly the
Jacobian, and the mass matrix that multiplies $\dot{\vect q}$ in the weak form is
the identity only in the $s$ measure.  All statements below are therefore made
with $\dd s$, and the derivative in $S_{ij}$ is $\dd_s=\dd/\dd s$; in the $\zeta$
coordinate $\dd_s=2\dd/\dd\zeta$, which is the factor of two that appears in
every strain-gradient entry.

The stiffness matrix is $S_{ij}=\int b_i'b_j'\dd s$.  For $i=j$ the chain rule
gives $b_j'(s)=2\sqrt{2j+1}\,P_j'(2s)$, so
\begin{equation}
\label{eq:app_S}
S_{jj}=4(2j+1)\int_{-1/2}^{1/2}\bigl(P_j'(2s)\bigr)^{2}\dd s
=2(2j+1)\,j(j+1),
\end{equation}
where $\int_{-1}^{1}(P_j')^{2}\dd\zeta=j(j+1)$ is the classical companion of the
Legendre equation and the factor $2$ collects $4\times\tfrac12$.  The weighted
form of that identity is the one that follows directly from the equation, and it
is worth deriving because it is the step usually skipped.  Multiplying
$\bigl[(1-\zeta^{2})P_j'\bigr]'=-j(j+1)P_j$ by $P_k$ and integrating by parts, the
boundary term carrying the factor $(1-\zeta^{2})$ vanishes at both ends and
\begin{equation}
\label{eq:app_S_weighted}
\int_{-1}^{1}(1-\zeta^{2})P_j'P_k'\,\dd\zeta
=\frac{2j(j+1)}{2j+1}\,\delta_{jk},
\end{equation}
which is the orthogonality of the derivatives under the weight
$(1-\zeta^{2})$; \eqref{eq:app_S} is its unweighted companion and is used here
only through the closed form it implies.  The discrete matrix checks both, and so
does an independent quadrature rule of a different order
(\cref{app:repro}): the assembled $S_{jj}$ equals $2j(j+1)(2j+1)$ for every
$j\le N-1$ at every degree used in this article, to $2\times10^{-11}$ or better
in absolute value.

For $i\neq j$ the parity of $P_i'P_j'$ leaves $S_{ij}\neq0$ only when $i$ and $j$
have the same parity, and
\begin{equation}
\label{eq:app_S_off}
S_{ij}=2\sqrt{(2i+1)(2j+1)}\int_{-1}^{1}P_i'P_j'\,\dd\zeta ,
\end{equation}
which is why \eqref{eq:app_S_matrix} is banded by two.  At $N=6$ the matrix is
\begin{equation}
\label{eq:app_S_matrix}
\mathsf S=
\begin{pmatrix}
0&0&0&0&0&0\\
0&12&0&18.330&0&22.978\\
0&0&60&0&80.498&0\\
0&18.330&0&168&0&210.60\\
0&0&80.498&0&360&0\\
0&22.978&0&210.60&0&660
\end{pmatrix},
\end{equation}
whose diagonal is \eqref{eq:app_S} entry by entry.  Equation
\eqref{eq:app_S_matrix} is the entire gradient-energy content of the
discretisation for a single field: it is real, symmetric and independent of the
material, and every strain-gradient block of \eqref{eq:block} is this matrix
multiplied by an elastic modulus and by $\lambda^{2}$.

Two further element matrices are needed.  The potential space of the grounded
problem is spanned by
\begin{equation}
\label{eq:app_pot}
\psi_j(s)=b_j(s)-r_j\,b_{j+2}(s),\qquad
r_j=\sqrt{\frac{2j+1}{2j+5}},\qquad j=0,\dots,N_\varphi-1 ,
\end{equation}
which vanishes at $s=\pm\tfrac12$ for every $j$ because $P_j(\pm1)=(\pm1)^j$ and
$r_j$ is exactly the ratio that makes the two terms cancel; $N_\varphi=N-2$,
since $\psi_j$ needs $b_{j+2}$ up to $j=N_\varphi-1$, i.e.\ $b_{N-1}$.  The two
matrices that carry the charge are then distinguished by which index is the test
index.  Testing Gauss's law with $\psi_i$ against the carrier basis $b_j$ gives
\begin{equation}
\label{eq:app_charge}
\bigl(\mathsf C\bigr)_{ij}\equiv\int_{-1/2}^{1/2}\psi_i\,b_j\,\dd s
=\delta_{ij}-r_i\,\delta_{j,i+2},
\qquad i=0,\dots,N_\varphi-1,\quad j=0,\dots,N-1 ,
\end{equation}
that is a unit matrix minus one band above the diagonal, with the coefficient
$r_i$ carrying the *row* index.  The reciprocal matrix that appears in the
carrier equation is its transpose,
\begin{equation}
\label{eq:app_charge_T}
\bigl(\mathsf C^{\mathsf T}\bigr)_{ij}\equiv\int_{-1/2}^{1/2}b_i\,\psi_j\,\dd s
=\delta_{ij}-r_j\,\delta_{i,j+2},
\end{equation}
which has the band below the diagonal and the coefficient $r_j$ on the column
index.  Both appear in \eqref{eq:block}: $-\mathsf C$ in the potential row and
$-\mathsf C^{\mathsf T}$ in the carrier row, entering through the single term
$-\int W_i^{*}W_j$ whose test and trial functions share the potentials.  Writing
$r_j$ where $r_i$ belongs, or transposing the two, gives a matrix that differs
from the assembled one by $0.2981$ in the largest entry at degree four --- which
is why the orientation is stated here rather than left to the reader.  Finally
the dielectric matrix is
$\kappa_{33}\int\psi_i'\psi_j'\dd s+\kappa_{11}k_y^{2}\int\psi_i\psi_j\dd s$,
with
\begin{equation}
\label{eq:app_diel}
\int_{-1/2}^{1/2}\psi_i'\psi_j'\,\dd s
=S_{ij}-r_jS_{i,j+2}-r_iS_{i+2,j}+r_ir_jS_{i+2,j+2},
\end{equation}
in which the second and third terms are equal because $\mathsf S$ is symmetric,
and the first term is not $S_{ij}$ alone: the grounded basis removes a band of
the stiffness exactly as it removes a band of the mass.

\subsection{Five worked matrix elements}

The entries below are evaluated at the operating point $k_xd=0.7362$,
$k_yd=0.1298$, $\omega=7.6334\,c_s/d$ and at polynomial degree $p=4$, i.e.\
$N=p+1=5$ basis functions per field, $n_{\mathrm{u}}=3N=15$, $n_\varphi=N-2=3$,
$n_\mu=N=5$ and $n_{\mathrm{w}}=5p+3=23$ unknowns.  That is the smallest degree at
which the parity structure and the grounded bands are both visible; the reported
results use $p=18$ to $30$, where the same formulas hold with larger matrices.
All entries are in the scaled units of the solver --- stresses divided by
$c_{44}$, fields by $\sqrt{c_{44}\kappa_{33}}$, permittivities by
$\kappa_{33}$ --- so the numbers are dimensionless.

\begin{enumerate}
\item \emph{Compressibility (mass) entry}, carrier row and column:
\begin{equation}
\label{eq:app_w1}
-\int_{-1/2}^{1/2}b_ib_j\,\dd s=-\delta_{ij}
\qquad\text{so that}\qquad
\bigl(\mathsf A_{00}\bigr)_{\mu_0\mu_0}=-1 ,
\end{equation}
reproduced as $-1$ to $7\times10^{-15}$.
\item \emph{Transport entry.}  The carrier gradient enters the weak form without a
power of $k_x$ through its transverse part and with one power through its axial
part, so the $k_x^{0}$ block carries the transverse wavenumber and the $k_x^{2}$
block carries the axial one:
\begin{equation}
\label{eq:app_w2}
\bigl(\mathsf A_{00}\bigr)_{\mu_i\mu_j}\Big|_{\text{transport}}
=-\frac{\ii\alpha_{\mathrm d}}{\omega}\Bigl[k_y^{2}\,\delta_{ij}+S_{ij}\Bigr],
\qquad
\bigl(\mathsf A_{11}\bigr)_{\mu_i\mu_j}\Big|_{\text{transport}}
=-\frac{\ii\alpha_{\mathrm d}}{\omega}\,\delta_{ij} ,
\end{equation}
so that the polynomial operator reproduces
$-\ii\alpha_{\mathrm d}\bigl(k_x^{2}+k_y^{2}+S_{ij}\bigr)/\omega$.  With
$\alpha_{\mathrm d}/\omega=0.8283/7.6334=0.108512$ and $k_y^{2}=0.016848$ the
diagonal of \eqref{eq:app_w2} is
\begin{equation}
\label{eq:app_w2_values}
-\ii\,1.828\times10^{-3},\quad -\ii\,1.3039,\quad -\ii\,6.5124,\quad
-\ii\,18.2314 \qquad (j=0,\dots,3),
\end{equation}
for $S_{jj}=0,12,60,168$, and the $k_x^{2}$ block is
$-\ii\,0.108509$.  The two pieces together reproduce the assembled coefficient of
the polynomial to $3.6\times10^{-15}$.  The growth with $j$ is the discrete
statement that a higher Legendre mode stores more gradient energy, and therefore
that diffusion penalises thickness structure.
\item \emph{Carrier--potential entry.}  With the orientations fixed by
\eqref{eq:app_charge}--\eqref{eq:app_charge_T}, the potential row of
$\mathsf A_{00}$ carries $-\mathsf C$ and the carrier row
$-\mathsf C^{\mathsf T}$:
\begin{equation}
\label{eq:app_w3}
\bigl(\mathsf A_{00}\bigr)_{\varphi_i\mu_j}
=-Q\bigl(\delta_{ij}-r_i\,\delta_{j,i+2}\bigr),
\qquad
\bigl(\mathsf A_{00}\bigr)_{\mu_i\varphi_j}
=\bigl(\mathsf A_{00}\bigr)_{\varphi_j\mu_i},
\end{equation}
so the diagonal is $-Q=-0.074635$ and the band is $+Qr_i$, that is $+0.033378$,
$+0.048860$, $+0.055630$ and $+0.059538$ for $i=0,\dots,3$; the first of these is
$Q/\sqrt{5}$ because $r_0=1/\sqrt{5}$, and the sign is positive because the
band of \eqref{eq:app_charge} is negative and the block carries the overall minus
of $-\int W_i^{*}W_j$.  In the carrier row the same numbers appear at
$(i,i-2)$, below the diagonal, with the coefficient $r_{i-2}$ carrying the
*trial* (potential) index; the two matrices are transposes of one another, and
the assembled blocks reproduce \eqref{eq:app_w3} to $9\times10^{-16}$.
\item \emph{Gradient-stiffness entry}, out-of-plane row and column:
\begin{equation}
\label{eq:app_w4}
\lambda^{2}\frac{c_{33}}{c_{44}}\,S_{ij}
=0.0625\times5.1395\times S_{ij}
\;=\;3.85,\ 19.27,\ 54.00,\ 115.64 \quad (j=1,2,3,4),
\end{equation}
growing as $j^{3}$ because $S_{jj}\propto j^{3}$: this is the growth that makes
the discrete operator badly scaled before the row equilibration of
\cref{sec:numerics_edge}, and $115.64$ at $j=4$ is the largest of these entries at
this degree.
\item \emph{Inertia entry}:
\begin{equation}
\label{eq:app_w5}
-\omega^{2}\delta_{ij}=-58.27\,\delta_{ij}\qquad(\omega=7.6334).
\end{equation}
The ratio of the inertia entry to the $j=1$ gradient entry is
$58.27/3.85=15.1$.  It compares two matrix entries at one fixed frequency and
nothing more: the frequency of a thickness mode is an eigenvalue of the coupled
problem \eqref{eq:quartic}, not the ratio of two entries, and no mode frequency is
inferred from it here.
\end{enumerate}

\subsection{Construction of the blocks from the weak equations}

Each of the nine blocks $\mathsf A_{rs}$ is the sum of the eight bilinear forms
listed in \eqref{eq:block}.  Written out with the operators of
\eqref{eq:sector_blocks} and the integrals of Appendix~\ref{app:derivations}.1,
the block is
\begin{equation}
\label{eq:app_block}
\begin{aligned}
\mathsf A_{rs}={}&\underbrace{\bigl\langle\mathsf B_r,\mathsf C\mathsf B_s
\bigr\rangle}_{\text{elastic}}
+\underbrace{\lambda^{2}\bigl\langle\partial_\zeta\mathsf B_r,
\mathsf C\,\partial_\zeta\mathsf B_s\bigr\rangle}_{\text{gradient}}
+\underbrace{\bigl\langle\mathsf G_r,\mathsf e\mathsf B_s\bigr\rangle
+\bigl\langle\mathsf B_r,\mathsf e^{\mathsf T}\mathsf G_s\bigr\rangle}_{\text{piezoelectric}}\\
&-\underbrace{\bigl\langle\mathsf G_r,\boldsymbol\kappa\mathsf G_s\bigr\rangle
}_{\text{dielectric}}
+\underbrace{f\bigl[\bigl\langle\mathsf G_r^{(2)},\partial_\zeta\mathsf B_s^{(2)}
\bigr\rangle+\bigl\langle\partial_\zeta\mathsf B_r^{(2)},\mathsf G_s^{(2)}
\bigr\rangle\bigr]}_{\text{flexoelectric}}
-\underbrace{\frac{\ii\alpha_{\mathrm d}}{\omega}
\bigl\langle\mathsf M\mathsf G_r,\mathsf M\mathsf G_s\bigr\rangle}_{\text{transport}}\\
&-\underbrace{\bigl\langle\mathsf W_r,\mathsf W_s\bigr\rangle}_{\text{carrier}}
-\underbrace{\omega^{2}\bigl\langle\mathsf U_r,\mathsf U\bigr\rangle\delta_{r0}
\delta_{s0}}_{\text{inertia}},
\end{aligned}
\end{equation}
with $\langle\vect a,\vect b\rangle=\int_{-1/2}^{1/2}\vect a^{\mathsf H}\vect b\,
\dd s$ in the measure of \eqref{eq:app_measures}.  The strain operator carries the $k_x$ power in its first slot,
$\mathsf B_0$ contains $\ii k_xb_j$ for $\varepsilon_{xx}$ and $\mathsf B_1$ contains $b_j$ for the
$\partial_x^1$ part, so that the $k_x$ power of a block is the sum of the powers
in its two factors; this is the bookkeeping that produces
\eqref{eq:polynomial} from \eqref{eq:app_block} without any further algebra.

\subsection{The edge-matching system}

At the edge $x=0$ the conditions \eqref{eq:bc_edge} are imposed weakly, one row
block per condition and each on the test functions of the field it multiplies:
the traction $\Sigma_{3j}$ on the $n_{\mathrm u}=3N$ displacement test functions,
the double traction $\tau_{3j3}$ on the same $n_{\mathrm u}$ test functions (a
block that exists only when the strain-gradient length is non-zero, since
$\tau_{3j3}$ vanishes with it), the normal electric displacement $D_z$ on the
$n_\varphi$ potential test functions and the normal carrier flux $J_z$ on the
$n_\mu$ carrier test functions.  With the edge ansatz \eqref{eq:edge_ansatz},
\begin{equation}
\label{eq:app_edge}
\vect y(0)=\vect y_{\text{inc}}+\sum_{j=1}^{M}a_j\,\vect y^{(j)},
\qquad
\mathsf B\vect a=-\vect b_{\text{inc}},
\end{equation}
where the matrix of the system and its right-hand side are the four row blocks
of the boundary operator $\mathsf R$ applied to the modes,
\begin{equation}
\label{eq:app_edge_R}
\mathsf R=\begin{pmatrix}
\mathsf R_{\Sigma}\\
\mathsf R_{\tau}\\
\mathsf R_{D}\\
\mathsf R_{J}
\end{pmatrix}\in\mathbb{C}^{m\times M},
\qquad
\mathsf B_{ij}=\bigl(\mathsf R\vect y^{(j)}\bigr)_i,\quad
\vect b_{\text{inc}}=\mathsf R\,\vect y_{\text{inc}},
\end{equation}
so that $m=2n_{\mathrm u}+n_\varphi+n_\mu$.  At degree four
($n_{\mathrm u}=15$, $n_\varphi=3$, $n_\mu=5$) this is $m=38$ rows against
$M=n_{\mathrm w}=23$ columns: the edge system is overdetermined by the fifteen
double-traction rows and is solved in the least-squares sense by a singular-value
decomposition.  With the flexoelectric coupling switched off the
$\mathsf R_\tau$ block is absent and $m=n_{\mathrm u}+n_\varphi+n_\mu=23=n_{\mathrm w}$
is square, which is the only case in which it is.  The system solved is the
column-scaled form of \eqref{eq:app_edge_R}, scaled per column by the norm of
that column and per row by the norm of that row; the rank condition and the
residual are quoted in \cref{sec:numerics_edge}, and the conditioning at the
operating point is $1.6\times10^{2}$ after scaling.

\subsection{Normalisation of the modal power}

The power \eqref{eq:Px_int} is bilinear in a mode and its conjugate, so the power
carried by the $j$-th outgoing mode scales with the square of its excitation
amplitude and with the normalisation of its eigenvector.  The eigenvector is
normalised by $\vect y^{(j)}\to\vect y^{(j)}/\sqrt{\mathsf y^{\mathsf H}\mathsf
M\mathsf y}$, which fixes its mass norm but leaves the *sign* of the power free:
the physical power is
\begin{equation}
\label{eq:app_power}
P_j=|a_j|^{2}\,\mathcal{P}_x\bigl(\vect y^{(j)}\bigr),
\end{equation}
and the degree-invariance of $P_j$ is the test that the normalisation has been
applied correctly.  At the operating point, for the backward mode,
\begin{equation}
\label{eq:app_power_numbers}
\begin{aligned}
n=18:&\quad \mathcal{P}_x=-0.0178923,\quad |a|=1.13749,\quad P_B=-0.0231481,\\
n=22:&\quad \mathcal{P}_x=-0.0154248,\quad |a|=1.22510,\quad P_B=-0.0231481,\\
n=26:&\quad \mathcal{P}_x=-0.0132455,\quad |a|=1.32205,\quad P_B=-0.0231481,\\
n=30:&\quad \mathcal{P}_x=-0.0113777,\quad |a|=1.42644,\quad P_B=-0.0231481,
\end{aligned}
\end{equation}
in which the raw power changes by \SI{57}{\percent} across the four degrees while
the physical power is constant to the seven digits shown.  That constancy is the
numerical statement of \eqref{eq:app_power} and it is the reason the reported
share is degree-independent at the $10^{-7}$ level of
\cref{tab:refinement}.

\subsection{One end-to-end evaluation of the reported share}

The reported number at the operating point is assembled from five quantities,
each of which is stored with every row of the data: the raw backward power
$\mathcal{P}_x$, the scattering amplitude $|a_B|$, the incident power at the
reference plane, the two attenuation exponents and the reference distance
$L=10\,d$.  With the values of \eqref{eq:app_power_numbers} at $n=22$,
\begin{equation}
\label{eq:app_RB}
R_B=\frac{|a_B|^{2}\bigl|\mathcal{P}_x\bigr|\,e^{2\Imag k_{Bx}L}}
{P_{\text{inc}}\,e^{2\Imag k_{ix}L}}
=\frac{0.0231505\times0.9998954}{1.0001081}
=\frac{0.0231481}{1.0001081}=0.0231456 ,
\end{equation}
where $|a_B|^{2}|\mathcal{P}_x|=1.2938\times0.0178923=0.0231505$ is the power at
the edge, the factor $e^{2\Imag k_{Bx}L}=0.9998954$ carries it to $x=-L$ with
$\Imag k_{Bx}=-5.2306\times10^{-6}$ per unit $d$, and the incident mode carries a
factor $e^{2\Imag k_{ix}L}=1.0001081$ with $\Imag k_{ix}=+5.4027\times10^{-6}$.
What the number measures
is the ratio of two physical powers, $-P_B/P_{\text{inc}}=0.0231481$, and the
two exponential factors are there to state at which plane each power was taken.
Reading the same row the other way, the backward mode carries
$\SI{2.31}{\percent}$ of the incident power at $x=-10d$, and
$\SI{2.31}{\percent}\times e^{-2.13\times10^{-4}}$ at $x=-20d$: the decay of the
share along the plate is governed by the *difference* of the two attenuation
constants, not by either one of them.
